% Generated draft: requires mathematical and source review.
\documentclass[11pt]{article}
\usepackage[margin=25mm]{geometry}
\usepackage{fontspec}
\setmainfont{Latin Modern Roman}
\usepackage{amsmath,amssymb,mathtools,unicode-math}
\setmathfont{Latin Modern Math}
\usepackage{enumitem,needspace,xurl,hyperref}
\setlength{\parindent}{0pt}
\setlength{\parskip}{5pt}
\setlength{\emergencystretch}{3em}
\newcommand{\N}{\mathbb{N}}
\newcommand{\Z}{\mathbb{Z}}
\newcommand{\Q}{\mathbb{Q}}
\newcommand{\R}{\mathbb{R}}
\newcommand{\C}{\mathbb{C}}
\newcommand{\F}{\mathbb{F}}
\newcommand{\A}{\mathbb{A}}
\newcommand{\PP}{\mathbb{P}}
\newcommand{\E}{\mathbb{E}}
\newcommand{\eps}{\varepsilon}
\newcommand{\dd}{\,\mathrm{d}}
\begin{document}
\section*{UnsolvedMath ID: 30003238}
\textbf{Source problem code:} catalogue entry
\subsection{Definitions and mathematical statement}
For a conventional interpretation of the displayed expressions, let \(\alpha,\beta\in\R\), let \(p\) be a fixed prime, and let \(q\) tend to infinity through positive integers. These domain conventions and the intended quantification over \(p\) are not explicit in the supplied statement and require confirmation. Define
\[
\|x\|=\inf_{n\in\Z}|x-n|,\qquad
v_p(q)=\max\{k\in\Z_{\ge0}:p^k\mid q\},\qquad
|q|_p=p^{-v_p(q)}.
\]
Use natural logarithms and restrict the expressions below to \(q\ge3\); this restriction does not affect their limiting behavior. Put
\[
L(q)=q\log q\log\log q,\qquad M(q)=q\log q.
\]
For \(w\in\{L,M\}\), define
\[
S_w=\left\{(\alpha,\beta)\in\R^2:
\liminf_{\substack{q\to\infty\\q\in\N}}w(q)\|q\alpha\|\|q\beta\|>0\right\},
\]
\[
T_{w,p}=\left\{\alpha\in\R:
\liminf_{\substack{q\to\infty\\q\in\N}}w(q)|q|_p\|q\alpha\|>0\right\}.
\]
Here \(\dim_H\) denotes Hausdorff dimension in the usual Euclidean metric. The supplied record presents
\[
\dim_H S_L=2,\qquad \dim_H T_{L,p}=1
\]
as background results; these assertions are reproduced as supplied and have not been independently verified.

The logarithmic-loss question is whether the corresponding full-dimension conclusions hold after removing \(\log\log q\), namely whether
\[
\dim_H S_M=2,\qquad \dim_H T_{M,p}=1.
\]
The record also asks whether ``these sets'' are winning and what happens when they are intersected with a planar curve. The winning question requires specification of the game, its parameters, and whether it concerns the original sets, the reduced-factor sets, or both. For a planar curve \(C\subseteq\R^2\), the planar set has the well-defined intersection \(S_L\cap C\), and likewise \(S_M\cap C\) for the reduced-factor version. However, the record specifies neither the class of curves nor a particular intersection conclusion, such as a Hausdorff-dimension formula. Since \(T_{w,p}\subseteq\R\), its intersection with a planar curve also requires an embedding, parametrization, or other interpretation that is not supplied. These portions therefore remain requests for a precise formulation rather than fully specified conjectures.
\subsection{Short English statement}
The supplied record states that two multiplicative Diophantine approximation sets have full Hausdorff dimension when the weight is \(q\log q\log\log q\). It asks whether full dimension persists with weight \(q\log q\), whether the relevant sets are winning, and how they meet planar curves. The prime quantifier, winning game, curve hypotheses, and intended intersection conclusions need clarification. The supplied open status is dated curation, not an independently verified conclusion.
\subsection{Sources}
\begin{itemize}
\item Supplied local UnsolvedMath registry record, \texttt{data/unsolvedmath/problems.json}, permanent ID \(30003238\), title: Removing Logarithmic Losses in Multiplicative Approximation; exact source problem code: \texttt{catalogue entry}. No individual proposer is identified.
\item Supplied report citation: Arbeitsgemeinschaft: Diophantine Approximation, Fractal Geometry and Dynamics (2017), Oberwolfach Reports, DOI 10.4171/owr/2016/48, pp.\ 2749--2792. \url{https://doi.org/10.4171/owr/2016/48}. The record assigns proposed year 2016; the supplied report citation uses 2017.
\item Supplied literature source: Badziahin, \emph{On multiplicatively badly approximable numbers}. \url{https://arxiv.org/abs/1101.1855}.
\item Supplied literature source: Oberwolfach Report 48. \url{https://doi.org/10.4171/owr/2016/48}.
\item Literature assessment dated exactly \texttt{2026-08-21}: the supplied curation labels the problem open and reports that logarithmic-factor constructions give full-dimensional sets, while removal of the extra \(\log\log q\) factor, the relevant winning property, and planar-curve intersection questions remain open in the stated generality. This is supplied dated curation derived from the OWR material, not independently verified literature assessment; its status and cited claims require independent checking.
\end{itemize}
\end{document}
